\documentclass[11pt]{article}

\usepackage{amsmath,amssymb}
\usepackage{xurl}
\usepackage[hidelinks]{hyperref}
\setlength{\emergencystretch}{2em}

\input{command}

\title{Issue 1230: How the Self-Improvement SDP Is Used}
\author{}
\date{2026-05-05}

\begin{document}
\maketitle

\section{At a glance}

\begin{itemize}
  \item \textbf{Status: discharged in \ghpr{1708}.}  The former issue was the full
  \textbf{strong-duality and complementary-slackness producer} for the
  self-improvement SDP, not just the already-fixed submeasurement interface.
  The current formalization proves this producer in the canonical finite-
  dimensional SDP model and transports it to the paper-facing statement.
  \item \textbf{Difficulty: medium.}  The underlying SDP argument is standard
  finite-dimensional convex duality, but the formalization must choose a
  canonical block SDP whose optimal primal witness has no slack block without
  assuming the auxiliary dominance condition $I\preceq Z$.
  \item \textbf{Estimated weight: medium.}  The completed proof uses a
  substantial project-local canonical SDP layer, while the analytic inputs are
  ordinary finite-dimensional matrix-order and separation facts.
  \item \textbf{Traceability.}  The issue is tracked in \ghissue{1230}; the
  canonical strong-duality route was completed in \ghpr{1708}, the obsolete
  dominance wrappers were removed in \ghpr{2422}, and the reduced
  \path{SdpStatement}/\path{sdp} interface was retired in \ghpr{2430}.
  \item \textbf{Implemented route:} project-local canonicalization, compactness
  and separation for the canonical SDP image cone, equality of primal and dual
  optima, complementary slackness from weak duality, saturation of the auxiliary
  slack block, and the self-improvement adapter.
  \item \textbf{Mathlib/project split:} roughly \textbf{35\% Mathlib} and
  \textbf{65\% project-local}.  Mathlib supplies matrix order, traces,
  positive-semidefinite cones, convex cones, and separation tools; the project
  must build the paper-specific block SDP and helper interface.
  \item \textbf{Key Mathlib inputs:} positive semidefinite matrices, trace pairing,
  convex cones and dual cones, finite-dimensional linear maps, and hyperplane
  separation for closed/proper cones.%
  \footnote{Representative APIs live around
  \texttt{Mathlib.LinearAlgebra.Matrix.PosDef},
  \texttt{Mathlib.Analysis.Matrix.Order},
  \texttt{Mathlib.Analysis.Convex.Cone.Dual},
  \texttt{Mathlib.Analysis.Convex.Cone.InnerDual}, and
  \texttt{Mathlib.Geometry.Convex.Cone.*}.}
\end{itemize}

\section{Key theorem forms to develop}

The following theorem forms are the ones needed for the LDT self-improvement
argument and are now represented in Lean.  They are deliberately stronger than
the local issue~196 submeasurement fix.

\begin{enumerate}
  \item \textbf{PSD effect bounds.}  For every polynomial $g$, prove
  \[
    0\preceq A_g\preceq I,
    \qquad
    A_g=\E_{\bu} A^{\bu}_{g(\bu)}.
  \]
  This supplies dual feasibility of the strict witness $Z=2I$ and nonnegative
  objective gain when slack mass is completed.

  \item \textbf{Canonical block SDP equivalence.}  The paper primal over
  submeasurements is represented by the canonical PSD variable
  $X=\operatorname{diag}(T_{g_1},\ldots,T_{g_M},S)$ with
  \[
    \sum_gT_g+S=I,
    \qquad
    \langle C,X\rangle=\sum_g\tr(T_gA_g).
  \]
  Its dual is the paper dual $Z\succeq A_g$.

  \item \textbf{Strong duality.}  The formal proof uses the strict witnesses
  $T_g=(2M)^{-1}I$ and $Z=2I$ for feasibility and proves finite-dimensional SDP
  zero gap by separating the canonical primal image cone: there are optimal
  $X_*$ and $Z_*$ with equal primal and dual values.

  \item \textbf{Saturated slackness witness.}  Choose an optimal primal witness
  with \textbf{no slack block}, i.e. $\sum_gT_g=I$, and a dual optimum $Z$ such
  that
  \[
    Z\succeq A_g,
    \qquad
    T_g(Z-A_g)=0
    \quad\text{for all }g.
  \]
  Saturation is obtained by completing primal slack into an outcome whose
  objective coefficient is positive semidefinite, without decreasing the
  objective.  Slackness is then recovered from zero duality gap and
  positive-semidefinite dual slacks.  This route avoids the stronger and
  generally unavailable hypothesis $I\preceq Z$.

  \item \textbf{Self-improvement SDP output.}  Package the previous item as the
  exact data consumed by the helper proof:
  \[
    T\text{ is a measurement},
    \qquad
    Z\succeq0,
    \qquad
    Z\succeq A_g,
    \qquad
    T_gZ=T_gA_g.
  \]
  The dominance-carrying interfaces that ask for the extra field $I\preceq Z$
  are auxiliary conditional routes, not consequences of the SDP dual constraints
  and not hypotheses of the paper-facing theorem.  The faithful paper route uses
  the saturated-witness argument above.%
  \footnote{The current bridge isolating this interface issue is
  \doclink{MIPStarRE/LDT/SelfImprovement/Theorems/Results/SdpMatrixBridge.lean}.}
\end{enumerate}

\section{The source SDP}

The self-improvement lemma in \cite{Ji2020LowIndividualDegree} uses one SDP.%
\footnote{The SDP is in \path{references/ldt-paper/self_improvement.tex},
lines~62--88.  The strict-feasibility and complementary-slackness paragraph is
in lines~168--190.}
For $A_g=\E_{\bu}A^{\bu}_{g(\bu)}$, the primal and dual are
\begin{align}
  \sup_{\{T_g\}}&\quad \sum_g\tr(T_gA_g),&
  T_g&\succeq0,& \sum_gT_g&\preceq I,\tag{\texttt{eq:primal-objective}}\\
  \inf_Z&\quad \tr(Z),&
  Z&\succeq A_g&&\text{for every }g.\tag{\texttt{eq:dual-objective}}
\end{align}
The primal object is a \textbf{submeasurement}.  The SDP lemma then selects an
optimal pair with
\[
  \sum_gT_g=I,
  \qquad
  T_gZ=T_gA_g.
  \tag{\texttt{lem:sdp}, \texttt{eq:slater}}
\]

\section{The dominance condition is auxiliary}

The dual constraint $Z\succeq A_g$ does not imply $Z\succeq I$.  In the scalar
one-outcome case with $A_g=\frac12 I$, the trace-minimizing dual solution is
$Z=\frac12 I$, which is feasible and optimal but does not dominate the identity.
Thus the implication
\[
  S\succeq0,\qquad SZ=0,\qquad Z\succeq A_g\ \forall g
  \quad\Longrightarrow\quad S=0
\]
is not a valid consequence of dual feasibility alone.  It becomes valid after
adding the stronger hypothesis $I\preceq Z$, because then $Z\succeq0$ and
$S(Z-I)\succeq0$ after commuting the positive factors; from $SZ=0$ this gives
$0\preceq -S$, hence $S=0$.

The corresponding matrix-order assertion is therefore deliberately conditional:
it is valid under the additional hypothesis $I\preceq Z$, but not under dual
feasibility alone.  Earlier versions of the formalization packaged the
resulting SDP data in declarations whose names ended in \leanid{WithDominance};
those wrappers have now been removed from the live SDP interface.  Their former
existence did not prove that an arbitrary optimal dual witness selected by
Slater duality satisfies $I\preceq Z$.

\section{How the proof uses the optimum}

The paper defines
\[
  H^u_h=A^u_{h(u)}T_hA^u_{h(u)},
  \qquad
  H_h=\E_{\bu}H^{\bu}_h.
  \tag{\texttt{self\_improvement.tex:203--220}}
\]
The proof uses \textbf{exactly three SDP outputs}: submeasurement bounds make
$H^u$ and $H$ submeasurements; saturation collapses sums over $T_g$; and
slackness replaces $A_g$ by $Z$ after left multiplication by $T_g$.  For example,
\[
  \E_{\bu}\sum_h\bra{\psi}(T_hA^{\bu}_{h(\bu)})\ot I\ket{\psi}
  =\sum_h\bra{\psi}(T_hZ)\ot I\ket{\psi}.
\]
Dual feasibility $Z\succeq A_g$ then supplies the boundedness certificate carried
through the induction.%
\footnote{The two slackness substitutions are in
\path{references/ldt-paper/self_improvement.tex}, lines~397--400 and 580--584.
The boundedness conclusion is restated in
\path{references/ldt-paper/inductive_step.tex}, lines~277--284, 315--322,
477--484, and 542--549.}

\section{Blueprint and formalization}

The blueprint matches the paper: submeasurement primal, strict feasible witness,
then a saturated slackness-producing optimum.%
\footnote{See \path{blueprint/src/chapter/ch07_self_improvement.tex}, especially
\texttt{lem:sdp-uniform-feasible-witness},
\texttt{lem:sdp-matrix-feasible-bounds},
\texttt{lem:sdp-matrix-slackness-output}, and \texttt{lem:sdp}.}
Lean keeps the feasible primal boundary explicit as a submeasurement, while the
selected saturated optimum is exposed to the slackness-carrying statement as a
measurement.  The construction supplied by \ghpr{1708} is the \textbf{producer}
for this stronger statement.%
\footnote{Representative declarations are \leanid{sdpStrictPrimalWeight},
\leanid{sdpComplementarySlacknessEquation}, and
\leanid{averagedSandwichedPolynomialSubMeas} in
\doclink{MIPStarRE/LDT/SelfImprovement/Defs.lean}.  The active source-facing
interface is \leanid{SdpStatementWithSlackness} in
\doclink{MIPStarRE/LDT/SelfImprovement/Theorems/Statements.lean}; the former
reduced \leanid{SdpStatement} interface was retired in \ghpr{2430}.}

\section{Conclusion}

There is a genuine distinction between the unconditional SDP statement used in
the paper and the former conditional dominance bridges.  Those former wrappers
proved correct conditional statements: if a canonical optimal pair is supplied
with $I\preceq Z$, then complementary slackness forces the slack block to vanish
and the pair yields a complete measurement satisfying the displayed slackness
equations.  This does not by itself prove the unconditional conclusion of
\texttt{lem:sdp}; the unconditional conclusion is instead proved by the
canonical strong-duality and saturation route.

The faithful completion of the paper route is now represented by
\leanid{MIPStarRE.LDT.SelfImprovement.matrixSdpCanonicalStrongDuality} and
\leanid{MIPStarRE.LDT.SelfImprovement.sdp\_statement\_with\_slackness}: the
formalization proves strong duality with complementary slackness for the
canonical block SDP, produces a saturated optimal witness without assuming
$I\preceq Z$, and connects that witness to the self-improvement interface.
The reduced theorem \leanid{sdp}, its \leanid{SdpStatement} record, and the
wrapper \leanid{selfImprovementHelperConstruction} have been removed; the
remaining paper-facing route is the slackness-carrying statement
\leanid{SdpStatementWithSlackness}.

\bibliographystyle{alpha}
\bibliography{references}

\end{document}
